\documentclass[10pt,a4paper]{amsart}
\usepackage[margin=19mm]{geometry}
\usepackage{amsmath,amssymb,mathtools}
\usepackage[scr=rsfs,cal=euler]{mathalfa}
\usepackage{xfrac}
\usepackage[unicode,hidelinks,bookmarks=false]{hyperref}
\newcommand{\ie}{i.e.\ }
\newcommand{\cf}{cf.\ }
\newcommand{\resp}{respectively\ }
\newcommand{\Subsection}{\subsection}
\providecommand{\nobreakdash}{\nobreak\hbox{-}}
\setlength{\emergencystretch}{2em}
\multlinegap0pt
\usepackage{mdframed}
\usepackage{mdframed}
\usepackage{mdframed}
\usepackage{mdframed}
\usepackage{amssymb}
\usepackage{mdframed}
\usepackage{tikz}
\usepackage{algorithm}\usepackage{algpseudocode}
\usepackage{xcolor}
\newtheorem{proposition}{Proposition}
\newcommand{\RR}{\mathbb{R}}
\newcommand{\cA}{\mathcal{A}}
\newcommand{\cF}{\mathcal{F}}
\newcommand{\cP}{\mathcal{P}}
\newcommand{\cX}{\mathcal{X}}
\newcommand{\cY}{\mathcal{Y}} 
\title{Clipped Stochastic Gradient Tracking For\\ Locally Smooth Functions}
\author{Leilei Mei, Junyu Zhang}
\date{Source: arXiv:2605.17027v1 (CC BY 4.0)}
\begin{document}
\maketitle

\noindent\emph{Source and licence.} Clipped Stochastic Gradient Tracking For\\ Locally Smooth Functions. Version arXiv:2605.17027v1 (CC BY 4.0), distributed by arXiv under the Creative Commons Attribution 4.0 International licence, http://creativecommons.org/licenses/by/4.0/, as recorded in the arXiv licence metadata for this exact version and shown on the abstract page https://arxiv.org/abs/2605.17027v1. Full licence notice: \texttt{licenses/arxiv-2605.17027-CC-BY-4.0.txt}.\par\smallskip

\noindent\textit{Editorial scope.} Counted unit: the article's proposition proposition:closed (source0.tex lines 260--278). The article's complete original proof of it is delivered with it (source0.tex lines 279--333, one \texttt{proof} environment of 55 lines). No mathematics is written, repaired or re-derived: the statement and the proof are the article's own lines, in the article's own notation, and no retained line is substituted or re-flowed.  No further source-local context is needed: every cross-reference of every retained line resolves inside the two delivered spans.  Nothing was omitted from any retained span, so the package carries no omission record.  No retained line contains a citation, so the wrapper carries no bibliography and no bibliography entry.  No retained line contains a figure environment, an image-inclusion command or drawing code, so no image is delivered and the package declares no asset dependency.  Editorial formatting only, in addition: an AMS \texttt{amsart} preamble that restates the article's own declarations and macros used by the retained lines, namely the environments \texttt{proposition} on separate counters, together with the standard packages those lines need.  The retained environments print in this document as Proposition 1 in order of appearance, whereas the article prints the same items with the numbering of its own full document.  The complete native source of the article as distributed in the arXiv e-print for this exact version is bundled unchanged as \texttt{sources/200761/source0.tex}.
\begin{proposition}[Closedness of $\cF_{\text{ruc}}^d$]
\label{proposition:closed}
The RUC function class $\cF_{\emph{ruc}}^d$ is closed under affine composition, positive linear combination, maximum/minimum, and multiplication/division operations:\vspace{0.1cm} \\
\textbf{\emph{(1).}}
Let $F \in \cF_{\mathrm{ruc}}^d$ and let 
$\mathcal{A}:\mathbb{R}^{d'} \to \mathbb{R}^d$ be globally Lipschitz continuous. That is,   $\exists L>0$ s.t.
\[
    \|\mathcal{A}(x) - \mathcal{A}(y)\|
    \le L \|x-y\|,
    \qquad \forall x,y \in \mathbb{R}^{d'}.
\]
For $\cX \in \cP(\mathbb{R}^{d'})$ define 
$\mathcal{A}(\cX) := \{\mathcal{A}(x): x\in\cX\}$ and
$G(\cX) := F(\mathcal{A}(\cX))$.
Then $G \in \cF_{\mathrm{ruc}}^{d'}$.\\
\textbf{\emph{(2)}.}  If $F,G\in\cF_{\emph{ruc}}^{d}$, then $\alpha F+\beta G\in\cF_{\emph{ruc}}^{d}$ for any $\alpha,\beta>0$.\vspace{0.1cm}\\
\textbf{\emph{(3)}.} If $F,G\in\cF_{\emph{ruc}}^{d}$, then both $ \max\{F,G\}\in\cF_{\emph{ruc}}^{d}$ and $\min\{F,G\}\in\cF_{\emph{ruc}}^{d}$.\vspace{0.1cm}\\
\textbf{\emph{(4)}.} If  $F,G\in\cF_{\emph{ruc}}^{d}$, then both $F\cdot G\in\cF_{\emph{ruc}}^{d}$ and $F/G\in\cF_{\emph{ruc}}^{d}$.
\end{proposition}

\begin{proof}
First, for ease of notation, let us denote relative difference as 
$$d_{F}^{\rm rel}(\cX,\cY): = \left|1-\max\left\{\frac{F(\cX)}{F(\cY)},\frac{F(\cY)}{F(\cX)}\right\}\right|.$$
Then it is not hard to observe that $d_{F}^{\rm rel}(\cX,\cY)\leq \epsilon$ if and only if  $\max\left\{\frac{F(\cX)}{F(\cY)},\frac{F(\cY)}{F(\cX)}\right\}\leq 1+\epsilon$. First, let us  prove (1). By the $L$-Lipschitz continuity of $\mathcal{A}$, the induced set mapping is also
$L$-Lipschitz continuous under the Hausdorff distance:
\[
    d_H(\mathcal{A}(\cX),\mathcal{A}(\cY))
    \le L\cdot d_H(\cX,\cY),
    \qquad \forall \cX,\cY \in \mathcal{P}(\mathbb{R}^{d'}).
\]
Indeed, for any $x\in\cX$ we can find $y\in\cY$ with $\|x-y\|\le d_H(\cX,\cY)$,
and the Lipschitz property of $\mathcal{A}$ yields
$\operatorname{dist}(\mathcal{A}(x),\mathcal{A}(\cY))
\le L\cdot d_H(\cX,\cY)$, then taking maximum over all $x\in\mathcal{X}$ proves the first half of bound on Hausdorff distance between $\cA(\cX)$ and $\cA(\cY)$; the other half  follows by symmetry.
Next, fix an arbitrary $\epsilon>0$.
Because $F\in\cF_{\mathrm{ruc}}^d$, there exists $\delta>0$ such that
$d_H(\cX,\cY)\le\delta$ implies
$d_F^{\mathrm{rel}}(\cX,\cY)\le\epsilon$ for all
$\cX,\cY\in\mathcal{P}(\mathbb{R}^d)$.
Let $\delta' := \delta/L$.
Then for any $\cX,\cY\in\mathcal{P}(\mathbb{R}^{d'})$ with
$d_H(\cX,\cY)\le\delta'$, we have
\[
    d_H(\mathcal{A}(\cX),\mathcal{A}(\cY))
    \le L\, d_H(\cX,\cY)
    \le L\delta' = \delta,
\]
and hence, by definition that $G(\cX):=F(\mathcal{A}(\cX))$, we obtain
\[
    d_G^{\mathrm{rel}}(\cX,\cY)
    = d_F^{\mathrm{rel}}(\mathcal{A}(\cX),\mathcal{A}(\cY))
    \le \epsilon,
\]
which shows $G\in\cF_{\mathrm{ruc}}^{d'}$.  This completes the proof of (1).

Next we prove (2). Because $F,G\in\cF_{\rm ruc}^{d}$, for any $\epsilon>0$, there exist $\delta_F,\delta_G>0$ s.t. 
$d_{F}^{\rm rel}(\cX,\cY)\leq\epsilon$ for any $d_{\rm H}(\cX,\cY)\leq \delta_F$, and $d_{G}^{\rm rel}(\cX,\cY)\leq\epsilon$ for any $d_{\rm H}(\cX,\cY)\leq \delta_G$. Taking $\delta = \min\{\delta_F,\delta_G\}$, the above argument indicates that 
$F(\cX)\leq (1+\epsilon)F(\cY)$ and $G(\cX)\leq (1+\epsilon)G(\cY)$ simultaneously hold for any $d_{\rm H}(\cX,\cY)\leq \delta$, consequently
$$\frac{\alpha F(\cX)+\beta G(\cX)}{\alpha F(\cY)+\beta G(\cY)}\leq 1+\epsilon.$$
By the symmetry between $\cX$ and $\cY$,   repeating the above discussion with interchanged $\cX$ and $\cY$ gives 
$$\max\left\{\frac{\alpha F(\cX)+\beta G(\cX)}{\alpha F(\cY)+\beta G(\cY)},\frac{\alpha F(\cY)+\beta G(\cY)}{\alpha F(\cX)+\beta G(\cX)}\right\}\leq 1+\epsilon,$$
which implies an $\epsilon$-small relative difference: $d_{\alpha F+\beta G}^{\rm rel}(\cX,\cY)\leq\epsilon$. Hence we complete  the proof of (2) and conclude that $\alpha F+\beta G\in \cF_{\rm ruc}^{d}$.  

Next, we prove (3). For any $\epsilon>0$, let us inherit the choice of $\delta$ from the proof of (2). For any $d_{\rm H}(\cX,\cY)\leq\delta$, we assume without loss of generality that $F(\cX)\geq G(\cX)$. Then we have  
$$\frac{\max\{F(\cX),G(\cX)\}}{\max\{F(\cY),G(\cY)\}}=\frac{F(\cX)}{\max\{F(\cY),G(\cY)\}}\overset{(i)}{\leq}\frac{F(\cX)}{F(\cY)}\leq1+\epsilon,$$ where (i) is because $F(\cY)\leq\max\{F(\cX),G(\cX)\}$. Similar to the proof of (2), we can further use the symmetry between $\cX$ and $\cY$ to show $d_{\max\{F,G\}}^{\rm rel}(\cX,\cY)\leq\epsilon$ and $\max\{F,G\}\in\cF_{\rm ruc}^{d}$.

For the minimum function, we assume without loss of generality that $F(\cY)\geq G(\cY)$, then  
$$\frac{\min\{F(\cX),G(\cX)\}}{\min\{F(\cY),G(\cY)\}} = \frac{\min\{F(\cX),G(\cX)\}}{G(\cY)}\overset{(ii)}{\leq}\frac{G(\cX)}{G(\cY)}\leq 1+\epsilon,$$
where (ii) is due to $G(\cX)\geq\min\{F(\cX),G(\cX)\}$. Analogous to the maximum case, we can again use the symmetry argument to show $d_{\min\{F,G\}}^{\rm rel}(\cX,\cY)\leq\epsilon$ and hence $\min\{F,G\}\in\cF_{\rm ruc}^{d}$. 



Finally, we prove (4). Given any fixed $\epsilon'>0$, let us inherit the choice of $\delta$ from the proof of (2) with  $\epsilon$ satisfying 
$(1+\epsilon)^2 = 1+\epsilon'$. Then direct computation already implies that $F\cdot G\in\cF_{\rm ruc}^{d}$, and we omit this simple and sheer computation for succinctness. The division is also simple. Note that the symmetric definition of relative difference indicates that $d_{\min\{G\}}^{\rm rel}(\cX,\cY) \equiv d_{\min\{1/G\}}^{\rm rel}(\cX,\cY)$ for any $\cX',\cY'\in\mathcal{K}(\RR^{d'})$. Hence $G\in\cF_{\rm ruc}^{d}$ indicates that $1/G\in\cF_{\rm ruc}^{d}$. Then $F/G\in\cF_{\rm ruc}^{d}$ will hold by viewing the division between $F$ and $G$ as the multiplication between $F$ and $1/G$.
\end{proof}

\end{document}
